
% 3 frames, 5 minutes
\begin{frame}[plain]{}\leavevmode
{\color{blue!50!black}\LARGE\bfseries Statistical and computational efficiency\\[2mm]in sequential decision making\par}
\vspace{9mm}
\begin{tikzpicture}
\node[box,text width=4.6cm] (a) {Learning which actions to take};
\node[box,text width=4.6cm,right=12mm of a] (b) {Learning a model from incomplete observations};
\draw[arr] (a) -- (b);
\end{tikzpicture}
\gap Thesis defense
\end{frame}

\begin{frame}{Decisions change the data we get}
\begin{center}
\begin{tikzpicture}[node distance=14mm]
\node[box] (s) {System now};
\node[box,right=of s] (a) {Choose an action};
\node[box,right=of a] (n) {System later};
\draw[arr] (s)--(a);\draw[arr] (a)--(n);
\node[smallbox,below=13mm of a] (d) {Observations};
\draw[arr] (n.south) |- (d.east);
\draw[arr] (d.west) -| (s.south);
\end{tikzpicture}
\end{center}
\pause
\gap Trying an action costs data. Its consequences may appear much later.
\pause
\gap We may also observe only part of the system we want to control.
\end{frame}

\begin{frame}{Two questions for this talk}
\begin{columns}[T]
\column{.47\textwidth}
{\Large\color{blue}1. Learning a policy}\gap
How can regression and policy updates learn to act with limited interaction?
\gap\textbf{Actor--critic methods}\\
Work with Jianqing Fan and Yuting Wei.
\column{.47\textwidth}
\pause
{\Large\color{green}2. Estimating a model}\gap
How can we differentiate a likelihood estimated using particles?
\gap\textbf{Differentiable particle filtering}\\
Work with A. J. Abkemeier, Kunyang He, Giles Hooker, and Edward Ionides.
\end{columns}
\pause\gap\medskip
At the end: how these two problems meet in control with hidden states.
\end{frame}
